% ===========================================================================
% Poste 20 — Revenu disponible (orchestrateur)
% ===========================================================================
\poste{20}{Le revenu disponible (orchestrateur)}{RD}

\pdesc Le dernier poste n'est pas un programme, mais l'\emph{orchestrateur}
du moteur : à partir des seules entrées brutes (situation, revenus, âges,
enfants), il enchaîne les postes 1 à 19 dans l'ordre des dépendances
(section~\ref{sec:architecture}), reconstruit les bases de revenu
intermédiaires et agrège le tout.

\pobj Produire le revenu disponible et sa décomposition complète --- le
résultat que l'application affiche.

\pregle L'enchaînement, pas à pas :
\begin{enumerate}
  \item \emph{Cotisations sur le revenu} : RRQ, RQAP, AE (actifs) ; FSS
        (retraités) --- postes 1 à 4 ;
  \item \emph{Postes-feuilles} : Sécurité de la vieillesse (poste~17) et
        aide de dernier recours (poste~13, qui consomme les cotisations) ;
  \item \emph{Bases intermédiaires} : $\RFN$, $\AFNI$ et $\RAL$
        (équations~\ref{eq:rfn} à~\ref{eq:ral}) ;
  \item \emph{Prime RAMQ} (poste~5) sur le $\RFN$, avec exonération
        individuelle si l'adulte reçoit l'aide sociale ou un SRG d'au
        moins \pc{94} du maximum ;
  \item \emph{Impôts} (poste~19), qui reçoivent la prime RAMQ (crédit
        médical des couples) et la déduction fédérale pour frais de
        garde ;
  \item \emph{Transferts} : Allocation famille et fournitures scolaires,
        prime au travail, solidarité, allocation-logement, soutien aux
        aînés, crédit de frais de garde (sur le $\RFN$) ; ACE, crédit TPS,
        ACT (sur l'$\AFNI$) ;
  \item \emph{Agrégation} :
        \[
        RD = R - C + T_{QC} - I_{QC} + T_{CA} - I_{CA} - G
        \]
        soit l'équation~\ref{eq:rd}, arrondie au cent.
\end{enumerate}

L'orchestrateur expose, outre le revenu disponible : les six composantes de
l'équation, le \emph{détail des dix-neuf postes} (chaque montant
individuel) et les trois bases intermédiaires --- c'est ce détail que
l'application affiche et que la suite de tests confronte au calculateur
officiel, ligne par ligne.

\pparams Ce poste n'a pas de paramètres propres : il consomme ceux des
postes 1 à 19. L'application permet d'ailleurs de \emph{modifier} n'importe
lequel de ces paramètres (page \emph{Comparer des paramètres},
section~\ref{sec:manuel}) : le moteur accepte un jeu de paramètres complet
en entrée, les valeurs officielles n'étant que le préréglage par défaut.

% ============================================================================
% GÉNÉRÉ par scripts/exemples-guide.ts — NE PAS ÉDITER À LA MAIN.
% Régénérer : npm run exemples (chaque chiffre est vérifié contre le moteur).
% ============================================================================
\pexemples L'équation~\eqref{eq:rd} appliquée à trois ménages types (montants 2025
en dollars ; chaque composante est la somme des postes détaillés dans les
sections précédentes) :

\begin{center}
\small
\begin{tabular}{l r r r}
\toprule
 & M4 & M6 & M10 \\
\midrule
Revenu de travail ou de retraite & \num{50000} & \num{35000} & \num{20000} \\
Cotisations (postes 1--5) & $-$\num{4622} & $-$\num{2725.81} & $-$\num{762.70} \\
Transferts québécois (postes 6--13) & \num{908.60} & \num{7673.50} & \num{3144.92} \\
Impôt du Québec (poste 19) & $-$\num{3996.40} & $-$\num{1759.24} & $-$\num{100.26} \\
Transferts fédéraux (postes 14--18) & \num{332.30} & \num{8879} & \num{10424.99} \\
Impôt fédéral (poste 19) & $-$\num{3453.30} & $-$\num{0} & $-$\num{197.85} \\
Frais de garde payés (poste 6) & $-$\num{0} & $-$\num{2000} & $-$\num{0} \\
\midrule
\textbf{Revenu disponible} & \textbf{39169.20} & \textbf{45067.45} & \textbf{32509.10} \\
\bottomrule
\end{tabular}
\end{center}
\medskip\noindent À lire en colonne : par exemple, M6 (famille monoparentale,
\D{35000} de travail) \emph{ajoute} \D{7673.50}
de transferts québécois et \D{8879} de
transferts fédéraux à son revenu, pour un revenu disponible de
\D{45067.45} — supérieur à son revenu brut. À l'autre bout,
M4 (personne seule, \D{50000}) cède \D{12071.70}
en prélèvements et ne reçoit que \D{1240.90}
de transferts : revenu disponible de \D{39169.20}.


\prefs L'ensemble des références des postes 1 à 19 ; ministère des
Finances du Québec, calculateur officiel du revenu disponible (outil de
validation).
